\documentclass[11pt]{article}

\usepackage[margin=1in]{geometry}
\usepackage[T1]{fontenc}
\usepackage{amsmath}
\usepackage{amssymb}
\usepackage{hyperref}
\usepackage{listings}
\setlength{\emergencystretch}{2em}
\urlstyle{same}

\hypersetup{
    hidelinks
}

\lstset{
    basicstyle=\ttfamily\scriptsize,
    breaklines=true,
    breakatwhitespace=false,
    frame=single,
    columns=fullflexible,
    keepspaces=true,
    numbers=left,
    numberstyle=\tiny,
    numbersep=6pt,
    showstringspaces=false,
    tabsize=4
}

\title{
Legibility Optimization for a UD5 Manipulator:
Joint-Space Trajectory Optimization with
Forward-Kinematics-Based Goal Inference
}

\author{
HRI Legible Robot Motion Project
}

\date{Research report and repository integration: October 14, 2026}

\begin{document}

\maketitle

\section*{Executive Summary}

This document specifies a legibility-optimization module
placed after a conventional MoveIt motion planner and before
robot execution.

The baseline planner produces a feasible joint-space trajectory

\[
Q^0 = [q_0^0,\ldots,q_{N-1}^0].
\]

The legibility optimizer modifies the interior joint
configurations while preserving the start and target
configurations.

Human-visible legibility is evaluated in task space:

\[
x_i = FK(q_i).
\]

The optimizer therefore separates robot kinematics from the
human goal-inference model.

The proposed implementation combines the cost-based
goal-inference idea of Dragan and Srinivasa with the
configuration-space / task-space separation emphasized by
Bodden et al.

The optimizer output is a candidate trajectory, not an
automatically safe execution trajectory. The candidate must
subsequently undergo trajectory retiming, robot-specific collision/path and
joint-limit validation, and simulation before explicitly authorized hardware
execution. An optional geometric precheck may precede retiming; validation
of the final retimed trajectory remains mandatory.

\paragraph{Repository integration status.}
This manuscript preserves the supplied report's UD5 terminology.
Inspection of repository revision \texttt{18a1b66} finds documentation for
ROS~1 Noetic and a UR5e, not a confirmed UD5 kinematic model.
No equivalence between UD5, UR5, and UR5e is assumed.
The motion node and planner are TODO scaffolds; the separate recorded-pose
controller tool is unchanged. The new code is an offline experimental
package, outside the catkin launch and execution path. The installed MoveIt
version, robot description, planning group, goal frame, FK, collision checker,
and retimer still require lab verification. See the companion
\texttt{integration.md} and \texttt{provenance.md} for the inspected files
and the changes made while integrating the supplied source.

\section{Introduction}

Legible robot motion communicates a robot's intended goal
through the movement itself.

For a reaching task with multiple candidate targets, an
efficient trajectory may be ambiguous early in the motion.
A legibility optimizer intentionally modifies the trajectory so
that an observer can infer the true target earlier.

Dragan and Srinivasa introduced a probabilistic formulation of
legible motion and optimized trajectories while constraining
excessive deviation using a trust region
\cite{dragan2013}.

Bodden et al. later proposed a flexible constrained
optimization formulation for robot arms and explicitly
separated robot configuration space, task space, and the model
of human interpretation \cite{bodden2018}.

This separation motivates the present system.

The planner and robot communicate using joint-space
trajectories, but the legibility objective operates on
end-effector positions obtained from forward kinematics.

\section{System Architecture}

The system is

\[
\begin{gathered}
\text{MoveIt / RRT-Connect baseline}
\longrightarrow \text{standalone legibility optimizer} \\
\longrightarrow \text{retiming}
\longrightarrow \text{robot-specific validation} \\
\longrightarrow \text{simulation}
\longrightarrow \text{explicit execution}.
\end{gathered}
\]

The optimizer receives:

\begin{itemize}
    \item joint names,
    \item joint positions at all trajectory waypoints,
    \item timestamps,
    \item the intended target index,
    \item Cartesian positions of all candidate targets,
    \item a forward-kinematics callback,
    \item joint limits,
    \item hard per-joint local deviation limits,
    \item optional velocity and acceleration limits.
\end{itemize}

The optimizer returns:

\begin{itemize}
    \item candidate joint positions,
    \item nominal timestamps,
    \item baseline and candidate legibility scores,
    \item baseline and candidate motion costs,
    \item realized cost ratio,

    \item constraint-status diagnostics, numerical acceptance, and
          physical validity reported separately (unassessed in the core).
\end{itemize}

\subsection{System Flowchart}

The Mermaid source for the system diagram is:

\begin{lstlisting}
flowchart LR
     A[MoveIt / RRT-Connect baseline]
         --> B[RobotTrajectory]

     B --> C[Legibility optimizer]

     G[Fixed goal positions]
         --> C

     T[True trial target]
         --> C

     K[Local FK model]
         --> C

     C --> D[Optimized candidate]

     D --> E[Optional geometric precheck]

     E --> F[Trajectory retiming]

     F --> H[Validate retimed path and robot constraints]

     H --> I[RViz / simulation]

    I --> J[Explicit authorized execution]
\end{lstlisting}

\section{Trajectory Representation}

The robot path is represented as

\[
Q =
[q_0,q_1,\ldots,q_{N-1}],
\]

where

\[
q_i \in \mathbb{R}^{d}.
\]

The corresponding timestamps are

\[
T =
[t_0,t_1,\ldots,t_{N-1}].
\]

The optimizer fixes

\[
q_0 = q_0^0
\]

and

\[
q_{N-1} = q_{N-1}^0
\]

and changes only the interior configurations.

\section{Forward Kinematics}

Forward kinematics maps each joint configuration to a visible
end-effector position:

\[
x_i = FK(q_i),
\qquad
x_i \in \mathbb{R}^{3}.
\]

The software interface is

\begin{lstlisting}[language=Python]
def fk_position(q, joint_names):
    """
    q: joint configuration, shape (D,)
    joint_names: MoveIt joint order

    returns:
        xyz in the fixed-goal reference frame,
        shape (3,)
    """
    raise NotImplementedError("Connect verified robot FK")
\end{lstlisting}

This interface deliberately hides the robot-specific
kinematic implementation from the optimizer.

\section{Goal-Inference Model}

Let

\[
\mathcal G =
\{G_1,\ldots,G_M\}
\]

be the possible goals.

Let

\[
G^*
\]

be the true target.

A general cost-based observer has the form

\[
P(G\mid\xi_{0:i})
\propto
P(G)
\exp
\left(
-\beta[
C(\xi_{0:i})
+
V_G(x_i)
-
V_G(x_0)
]
\right).
\]

Because the prefix cost is common to every candidate goal, it
cancels during normalization.

For the first implementation, use

\[
V_G(x)
=
\|x-G\|_2.
\]

Then

\[
P(G\mid x_i)
=
\operatorname{softmax}_G
\left[
\beta
(
\|x_0-G\|
-
\|x_i-G\|
)
+
\log P(G)
\right].
\]

The parameter

\[
\beta > 0
\]

controls how strongly progress differences influence the
observer model.

The implementation uses the smooth distance
$d_\epsilon(x,G)=\sqrt{\|x-G\|_2^2+10^{-12}}$ in the same units as the
goals, and subtracts the maximum logit before exponentiation.
Priors must be finite, positive, and normalized. This progress-based
observer is Dragan-inspired; it is not a reproduction of the original
functional-gradient algorithm or of Bodden et al.'s complete objective.

\section{Legibility Objective}

Let

\[
p_i =
P(G^*\mid x_i).
\]

Normalize time:

\[
\tau_i =
\frac{t_i-t_0}
{t_{N-1}-t_0}.
\]

Use

\[
w_i =
(1-\tau_i)^\gamma.
\]

The legibility objective is

\[
L(Q)
=
\frac{
\sum_{i=1}^{N-2}
w_i p_i
}{
\sum_{i=1}^{N-2}w_i
}.
\]

The final waypoint is intentionally excluded because arriving at
the target should not dominate a metric intended to measure
pre-arrival inference.

The initial waypoint is also excluded because it contains no motion
evidence. This is a waypoint-weighted score, not a quadrature approximation
of a time integral. Nonuniform timestamps are supported, but changing
waypoint density may change the score. Comparisons should use the same
sampling rule and report any resampling after retiming.

\section{Joint-Space Cost Budget}

Define

\[
\Delta t_i =
t_i-t_{i-1}.
\]

The discrete squared-velocity cost is

\[
C_q(Q)
=
\sum_{i=1}^{N-1}
\frac{
\|q_i-q_{i-1}\|_2^2
}{
\Delta t_i
}.
\]

Let

\[
C_0 =
C_q(Q^0).
\]

The candidate must satisfy

\[
C_q(Q)
\le
\rho C_0,
\]

where

\[
\rho =
\texttt{cost\_ratio}.
\]

The value of $\rho$ should be experimentally calibrated rather
than copied blindly from prior work.

\section{Regularization}

To prevent unnecessary internal joint movement, define

\[
D(Q,Q^0)
=
\frac{1}{(N-2)d}
\sum_{i=1}^{N-2}
\|q_i-q_i^0\|_2^2.
\]

The objective minimized numerically is

\[
J(Q)
=
-L(Q)
+
\lambda_DD(Q,Q^0).
\]

\subsection{Hard Local Deviation Limits}

The soft penalty $D$ does not enforce a local bound. The integrated prototype
therefore additionally requires a finite positive vector $\delta\in\mathbb R^d$:
\[
 |q_{ij}-q^0_{ij}|\le\delta_j,
 \qquad i=1,\ldots,N-2,\quad j=1,\ldots,d.
\]
SLSQP bounds intersect these local intervals with the supplied joint limits.
Values use the corresponding joint's units (radians for revolute joints,
metres for prismatic joints); continuous joints require explicitly chosen
unwrapped intervals. No robot-specific default bounds are assumed.

\section{Complete Optimization Problem}

The proposed optimization is

\[
\begin{aligned}
Q^*
=
\arg\min_Q
\quad&
-L(Q)
+
\lambda_DD(Q,Q^0)
\\
\text{subject to}\quad
&
q_0=q_0^0,
\\
&
q_{N-1}=q_{N-1}^0,
\\
&
q^{\min}\le q_i\le q^{\max},
\\
&
|q_{ij}-q^0_{ij}|\le\delta_j,
\\
&
|\dot q_i|
\le
\dot q^{\max},
\\
&
|\ddot q_i|
\le
\ddot q^{\max},
\\
&
C_q(Q)
\le
\rho C_q(Q^0),
\\
&
Q\in\mathcal Q_{\mathrm{free}}.
\end{aligned}
\]

In the first implementation, membership of
$\mathcal Q_{\mathrm{free}}$ must be verified after numerical optimization
and retiming using a robot-specific MoveIt path/collision validator.
That validator is not implemented by this prototype: its interface raises
\texttt{NotImplementedError}. Linear densification is a diagnostic aid and
does not establish continuous collision freedom.

Optional velocity and acceleration bounds constrain finite differences:
\[
 v_i=\frac{q_{i+1}-q_i}{t_{i+1}-t_i},\qquad
 a_i=\frac{v_i-v_{i-1}}
 {\tfrac12[(t_{i+1}-t_i)+(t_i-t_{i-1})]}.
\]
These checks do not certify controller interpolation, jerk, endpoint
derivatives, or dynamic feasibility. Unconfigured dynamic checks are
reported as unassessed, not as successful. Retiming must rebuild derivatives;
if it changes positions or sampling, repeat geometric, task, local-deviation,
and motion-budget checks on an explicitly defined comparison grid.

Numerical acceptance requires solver success, finite correctly shaped
output, independently checked bounds and cost budget, and the configured
minimum improvement in $L$. The solver tolerance is not a physical safety
tolerance. The default postcheck tolerance is $10^{-7}$ in joint/dynamic
units; the cost tolerance is $10^{-7}\max(1,\rho C_0)$.
Invalid solver output is rejected even if a diagnostic baseline fallback
would otherwise satisfy all inequalities.

\section{Implementation}

The complete optimizer used by this document is embedded below.

\lstinputlisting[
    language=Python,
    caption={optimizer\_core.py}
]{../../experimental/legibility_optimizer/optimizer_core.py}

\section{ROS Adapter}

The adapter is intentionally separated from the optimizer. Repository
evidence supports a ROS~1 message interface, but not a concrete installed
MoveIt version or robot-specific kinematic API. The report's publishing node
is adapted into a read-only function taking a baseline, goal set, and target
from one trial together. It retains joint names and nominal seconds, checks
the explicit FK/goal frame labels, rejects unsupported multi-DOF paths, and
returns a diagnostic candidate. It creates no publishers or controllers and
does not reuse stale velocity, acceleration, or effort fields.

\lstinputlisting[
    language=Python,
    caption={ros1\_adapter.py: offline trajectory extraction and optimization}
]{../../experimental/legibility_optimizer/ros1_adapter.py}

\lstinputlisting[
    language=Python,
    caption={robot\_interfaces.py: missing robot integrations fail explicitly}
]{../../experimental/legibility_optimizer/robot_interfaces.py}

\section{Unit Tests}

\lstinputlisting[
    language=Python,
    caption={test\_optimizer\_core.py}
]{../../tests/experimental/test_optimizer_core.py}

\lstinputlisting[language=Python,caption={ROS message adapter tests}]
{../../tests/experimental/test_ros1_adapter.py}

\lstinputlisting[language=Python,caption={Execution isolation tests}]
{../../tests/experimental/test_no_execution.py}

\section{Verification Plot}

The primary numerical diagnostic is a plot of

\[
P(G^*\mid\xi_{0:t})
\]

for the baseline and optimized trajectory.

\begin{lstlisting}[language=Python]
plt.plot(
    time_baseline,
    probability_baseline,
    label="MoveIt baseline",
)

plt.plot(
    time_candidate,
    probability_candidate,
    label="Legibility candidate",
)

plt.xlabel("Normalized trajectory time")

plt.ylabel("P(true goal)")
plt.ylim(0.0, 1.0)
plt.legend()
plt.show()
\end{lstlisting}

The runnable companion demo uses a synthetic identity FK, writes CSV/NPZ/JSON
diagnostics, and optionally saves the probability and path plots. Its numerical
scores are not evidence of robot validity or human intent recognition.
\lstinputlisting[language=Python,caption={Offline demo and verification plot}]
{../../experimental/legibility_optimizer/demo.py}

\section{MoveIt Validation}

An optimization result must not be interpreted as executable
merely because the numerical solver reports success.

The candidate must subsequently be checked for:

\begin{enumerate}
    \item preserved start configuration;
    \item preserved target configuration;
    \item joint-position limits;
    \item velocity and acceleration limits;
    \item self collision;
    \item environmental collision;
    \item task constraints;
    \item valid trajectory timing;
    \item controller compatibility.
\end{enumerate}

The trajectory should be tested in simulation before physical
execution.

\section{Repository Reuse}

The report's highest-priority conceptual reference implementation is

\begin{verbatim}
https://github.com/uwgraphics/trajectoryoptimizer-public
\end{verbatim}

Important locations are

\begin{verbatim}
opt_example.py
trajopt/robot/builtinRobots.py
trajopt/spacetime/builtinObjectives.py
trajopt/spacetime/objective.py
trajopt/spacetime/builtinConstraints.py
trajopt/spacetime/constraint.py
trajopt/solver/slsqp.py
trajopt/solver/solverInterface.py
docs/TR1833.pdf

\end{verbatim}

Its conceptual separation among robot model, objectives,
constraints, and solver is useful, but the old Python-2
implementation should not automatically become a dependency of
the new ROS system. No code from these third-party repositories is copied
or vendored in this integration. Any future reuse requires checking the
license of the exact upstream revision and retaining required notices.

Other investigated repositories are:

\begin{verbatim}
https://github.com/khoilsr/hrc_legible_motion_generation

https://github.com/shiwenli0908/Legibility-Modulator

https://github.com/isir-socialrobotics/legibot

https://github.com/munawarkazmi/legibility-bounds
\end{verbatim}

\section{Experimental Parameters}

The first parameters to sweep are

\[
\rho
\in
\{1.00,1.05,1.10,1.20,1.30\}
\]

and

\[
\gamma
\in
\{0,1\}.
\]

The parameter $\beta$ should also be evaluated because it
controls the confidence of the model observer.

These are engineering test ranges and not claims about optimal
human perception.

\section{Lab Validation Checklist}

Before physical execution:

\begin{enumerate}

    \item load a saved MoveIt baseline;
    \item verify joint-name ordering;
    \item verify timestamps;
    \item verify FK against MoveIt/RViz;
    \item verify goal positions and FK use the same frame;
    \item optimize one target;
    \item inspect the end-effector path;
    \item inspect the legibility probability curve;
    \item check the cost ratio;
    \item run MoveIt collision/path validation;
    \item retime the trajectory;
    \item revalidate the retimed trajectory;
    \item execute in simulation;
    \item execute on hardware at conservative speed;
    \item record commanded and measured joint states.
\end{enumerate}

\section{Implementation Timeline}

The following preserves the report's proposed schedule, not a record of
completed robot integration or experiments.

\begin{lstlisting}
gantt
    title Legibility Optimizer Integration
    dateFormat YYYY-MM-DD

    section Core
    Mathematical formulation
        :done, math, 2026-10-14, 1d

    Core optimizer + tests
        :core, after math, 1d

    section Robot Integration
    Wire actual FK
        :fk, after core, 1d

    Parse MoveIt trajectory
        :io, after core, 1d

    MoveIt path validation
        :valid, after fk, 1d

    Trajectory retiming
        :retime, after valid, 1d

    section Evaluation
    RViz / simulation tests
        :sim, after retime, 1d

    cost_ratio sweep
        :sweep, after sim, 1d

    first physical tests
        :robot, after sweep, 1d

    participant pilot
        :pilot, after robot, 2d
\end{lstlisting}

\section{Limitations}

The mathematical observer is a model of human inference, not a
measurement of human inference.

A larger value of $L(Q)$ therefore means only that the candidate
is more legible according to the selected observer model.

Human-subject evaluation remains necessary to determine whether
the modified robot motion actually allows people to identify the
target earlier or more accurately.

The first optimizer also uses only end-effector position.
Future work may incorporate end-effector orientation, velocity,
observer viewpoint, learned perception models, or alternative
intent-expressiveness formulations.

\section{Conclusion}

The key architectural decision is to optimize the robot in
joint space while evaluating communicative intent through
forward kinematics in task space.

This keeps the trajectory physically meaningful to the robot
while allowing the legibility objective to reason about the
motion that a human observer actually sees.

The resulting Legibility Optimizer remains a post-planning
module:

\[
\text{baseline planning}
\rightarrow
\text{legibility optimization}
\rightarrow
\text{retiming}
\]
\[
\longrightarrow
\text{robot validation}
\rightarrow
\text{explicit execution}.
  \]

  \begin{thebibliography}{9}

  \bibitem{dragan2013}
  A. Dragan and S. Srinivasa,
  ``Generating Legible Motion,''
  Robotics: Science and Systems,
  2013.
  doi:10.15607/RSS.2013.IX.024.

  \bibitem{bodden2018}
  C. Bodden, D. Rakita, B. Mutlu, and M. Gleicher,
  ``A Flexible Optimization-Based Method for Synthesizing
  Intent-Expressive Robot Arm Motion,''
  International Journal of Robotics Research,
  vol. 37, no. 11,
  2018.
  doi:10.1177/0278364918792295.

  \bibitem{amirian2024}
  J. Amirian, M. Abrini, and M. Chetouani,
  ``Legibot: Generating Legible Motions for Service Robots
  Using Cost-Based Local Planners,''
  arXiv:2404.05100,
  2024.

  \end{thebibliography}

  \end{document}
